% Shared macros for the layered model figures (m1_layers, m2_layers):
% glossy spheres, a tilted mechanism plane and key-variable chain boxes.
% Shared colours for the TikZ schematics; same roles and hex values as
% Code/figures_matlab/paper_style.m (mycolor.space, base #1E5F74).
%   water (P, I, M, W-bar)  teal   | vegetation (V, r, d)  green
%   wind and sand (E, S)    brown  | highlight             orange
%   limits and caps         red dashed
\definecolor{pWater}{HTML}{1E5F74}
\definecolor{pWaterMid}{HTML}{659FB6}
\definecolor{pWaterLight}{HTML}{DCEBF1}
\definecolor{pTeal}{HTML}{00828C}
\definecolor{pVeg}{HTML}{2B633C}
\definecolor{pVegMid}{HTML}{4AC087}
\definecolor{pVegLight}{HTML}{E3F3E8}
\definecolor{pSand}{HTML}{8F7159}
\definecolor{pSandLight}{HTML}{FFEACE}
\definecolor{pAccent}{HTML}{F49431}
\definecolor{pLimit}{HTML}{D03B3B}
\definecolor{pInk}{HTML}{1F2B33}
\definecolor{pAxis}{HTML}{344A53}
\definecolor{pGrey}{HTML}{97AFB9}
\definecolor{pGreyDark}{HTML}{647B84}
\definecolor{pGrid}{HTML}{E6EBED}

\makeatletter
% Lighter version of TikZ's ball shading: white highlight, no black rim.
\pgfdeclareradialshading[tikz@ball]{glossy}{\pgfqpoint{-9bp}{10bp}}{%
  color(0bp)=(white); color(6bp)=(tikz@ball!18!white); color(19bp)=(tikz@ball!48!white);
  color(31bp)=(tikz@ball!78!white); color(50bp)=(tikz@ball)}
\tikzoption{glossy color}{\pgfutil@colorlet{tikz@ball}{#1}\def\tikz@shading{glossy}\tikz@addmode{\tikz@mode@shadetrue}}
\makeatother
\tikzset{
  tie/.style={draw=pAxis!75, line width=0.6pt},
  cbox/.style={draw=#1, rounded corners=2pt, line width=0.8pt, fill=white, text=#1,
    minimum width=1.45cm, minimum height=0.70cm, align=center, font=\normalsize},
  lk/.style={-{Stealth[length=5.5pt]}, line width=1pt, draw=#1},
  num/.style={circle, fill=white, draw=#1, text=#1, inner sep=0.4pt,
    font=\scriptsize\bfseries, minimum size=11pt},
  balltext/.style={align=center, font=\footnotesize, text=pInk, inner sep=0pt},
  law/.style={font=\footnotesize, text=pInk, align=left},
}
% \ball{x}{y}{r}{colour}{text}
\newcommand{\ball}[5]{%
  \shade[glossy color=#4!62] (#1,#2) circle (#3);
  \draw[#4!70, line width=0.45pt] (#1,#2) circle (#3);
  \node[balltext, text width=1.95*#3 cm] at (#1,#2) {#5};}
% \splitball{x}{y}{r}{left colour}{right colour}
\newcommand{\splitball}[5]{%
  \begin{scope}
    \clip (#1-#3-0.1,#2-#3-0.1) rectangle (#1-0.03,#2+#3+0.1);
    \shade[glossy color=#4!80] (#1-0.03,#2) circle (#3);
  \end{scope}
  \begin{scope}
    \clip (#1+0.03,#2-#3-0.1) rectangle (#1+#3+0.1,#2+#3+0.1);
    \shade[glossy color=#5!80] (#1+0.03,#2) circle (#3);
  \end{scope}}
% \plane{x0}{y0}{width}{skew}{depth}{colour}: opaque base of the tilted plane
\newcommand{\plane}[6]{%
  \fill[#6!30] (#1,#2) -- ++(#3,0) -- ++(#4,#5) -- ++(-#3,0) -- cycle;}
% translucent top coat so lines passing beneath the plane stay faintly visible
\newcommand{\replane}[6]{%
  \shade[left color=#6!50, right color=pWater!38, opacity=0.70]
    (#1,#2) -- ++(#3,0) -- ++(#4,#5) -- ++(-#3,0) -- cycle;
  \draw[#6!80!black, line width=0.5pt] (#1,#2) -- ++(#3,0) -- ++(#4,#5) -- ++(-#3,0) -- cycle;}
\newcommand{\factor}[3]{\node[font=\footnotesize\itshape, text=pInk!85] at (#1,#2) {#3};}
\newcommand{\uparrowshape}[3]{% x, y, colour
  \node[single arrow, shading=axis, top color=#3!20, bottom color=#3!80, draw=none,
    rotate=90, minimum height=0.55cm, minimum width=0.50cm, single arrow head extend=0.10cm,
    inner sep=0pt] at (#1,#2) {};}
\newcommand{\layerlabel}[3]{\node[font=\small\bfseries\sffamily, text=#2, anchor=west] at (14.65,#1) {#3};}
